% !TeX root = ../../main.tex
% ------------------------------------------------------------------
%  4.2 Arquitectura física --- título e introducción
%  Fragmento del Subdocumento 4.2, traido desde contenido.tex con
%  \parte{partes/02_0_arquitectura_fisica}. Las rutas son relativas a la
%  carpeta del subdocumento, nunca a la de main.tex.
%
%  NUMERACION. El índice obligatorio (aclaraciones, sección 11) numera este
%  contenido como 4.2 y sus subtítulos como 4.2.X. El título se imprime como
%  "4.2 Arquitectura física", sin el rótulo "Capítulo" (formato definido en
%  contenido.tex). La antigua "Visión general" es ahora esta introducción.
% ------------------------------------------------------------------

% Tablas sin palabras cortadas (aclaraciones, sección 5: "sin palabras
% cortadas por columnas angostas"). Rige desde aquí hasta el final del
% subdocumento, incluido el anexo T-11.
\ifdefined\LFXcuerpotablaorig\else
  \let\LFXcuerpotablaorig\LFXcuerpotabla
  \renewcommand{\LFXcuerpotabla}{%
    \LFXcuerpotablaorig\hyphenpenalty=10000\exhyphenpenalty=10000\relax}
\fi

\chapter{Arquitectura física}
\label{cap:4-2-arquitectura-fisica}

La arquitectura física declara dónde vive cada componente de la solución, qué servicios se contratan en la nube, cómo se conectan los sitios entre sí y con la nube, qué ocurre cuando falla cada conexión y cuánta capacidad se provee.

Esa arquitectura materializa el despliegue híbrido obligatorio del Artículo 16\textdegree{} de las Bases Administrativas. La carga principal corre en nube pública AWS, con la región primaria en sa-east-1 (São Paulo, Brasil) y la recuperación ante desastres en us-east-1 (Norte de Virginia, EE. UU.), y los componentes on-premise garantizan la continuidad de la operación de bodega, plataformas y terreno durante los cortes del enlace. La solución cubre las seis instalaciones de la compañía ---los centros de distribución de Talca y de Concepción, las plataformas de cross-docking de Curicó, Chillán y Los Ángeles, y la casa matriz de Talca---, además de la calle y los 14.200 puntos de entrega. Cinco de esas instalaciones alojan cómputo; la casa matriz, contigua al centro de distribución de Talca, solo tiene oficinas y trabaja contra la nube. La séptima instalación, que el caso proyecta a tres años, se agrega por parametrización, sin cambios de desarrollo ni crecimiento de servidores en el CD Talca. Dos principios gobiernan el diseño:

\begin{itemize}
  \item \textbf{Carga principal en nube:} el núcleo transaccional de preventa, reparto y canal moderno, la analítica, la integración y el almacenamiento consolidado se ejecutan en AWS sa-east-1, con escalado elástico para el peak de septiembre.
  \item \textbf{Autonomía total de la operación de bodega:} la preparación, el despacho, la recepción y el conteo corren siempre contra los servidores del centro de distribución, no contra la nube, de modo que la conexión a Internet no interviene en la operación de bodega.
\end{itemize}

Este apartado se ordena en seis partes. La sección~\ref{sec:implementos} resume y analiza el equipamiento y el software que se proveen; su especificación elemento por elemento está en el Formulario T-11, que acompaña a este subdocumento como archivo propio (LAFROX-Formulario-T-11). La sección~\ref{sec:emplazamiento} justifica dónde vive cada componente, en la nube o en las instalaciones, conforme a los criterios del Artículo 16.2. La sección~\ref{sec:servicios-nube} presenta los servicios contratados en la plataforma de nube. La sección~\ref{sec:despliegue} describe cómo la solución se pone en marcha y se mantiene en operación: ambientes y ciclo de entrega, red de despliegue, alta disponibilidad, recuperación ante desastres, respaldos y su verificación. La sección~\ref{sec:conexiones} describe las conexiones, sus puntos de falla y la contingencia de cada uno. La sección~\ref{sec:dimensionamiento} deriva el dimensionamiento y el plan de capacidad desde la volumetría del caso. Las especificaciones del data center primario y del secundario se desarrollan en el apartado~\ref{cap:4-3-data-center}, a continuación de este.
